% SPDX-License-Identifier: MIT-0
\documentclass[11pt]{article}
\usepackage[T1]{fontenc}
\usepackage[utf8]{inputenc}
\usepackage{lmodern}
\usepackage[a4paper,margin=25mm,headheight=15pt]{geometry}
\usepackage{amsmath,amssymb,amsthm,mathtools}
\usepackage{microtype,booktabs,array,longtable}
\usepackage[dvipsnames]{xcolor}
\usepackage{enumitem,fancyhdr}
\usepackage[colorlinks=true,linkcolor=MidnightBlue,citecolor=MidnightBlue,urlcolor=MidnightBlue]{hyperref}
\hypersetup{pdftitle={A parity factorization and sharp stabilization for Cigler's Hankel polynomials},pdfauthor={AI-assisted research draft prepared for Vladimir Reshetnikov}}
\definecolor{Ink}{HTML}{16324F}
\newtheorem{theorem}{Theorem}[section]
\newtheorem{lemma}[theorem]{Lemma}
\newtheorem{proposition}[theorem]{Proposition}
\newtheorem{corollary}[theorem]{Corollary}
\theoremstyle{definition}
\newtheorem{definition}[theorem]{Definition}
\newtheorem{conjecture}[theorem]{Conjecture}
\theoremstyle{remark}
\newtheorem{remark}[theorem]{Remark}
\newcommand{\Z}{\mathbb Z}
\newcommand{\Q}{\mathbb Q}
\newcommand{\N}{\mathbb N}
\newcommand{\CT}{\operatorname{CT}}
\newcommand{\diag}{\operatorname{diag}}
\newcommand{\U}{\mathcal U}
\newcommand{\V}{\mathcal V}
\newcommand{\Ht}{\mathcal H}
\newcommand{\F}{\mathcal F}
\newcommand{\B}{\mathcal B}
\newcommand{\ord}{\operatorname{ord}}
\newcommand{\coeffge}{\succeq}
\newcommand{\Dlt}{\Delta}
\newcommand{\eps}{\varepsilon}
\newcommand{\bracket}[1]{\left[#1\right]}
\DeclareMathOperator{\sgn}{sgn}
\numberwithin{equation}{section}
\pagestyle{fancy}
\fancyhf{}
\fancyhead[L]{\small\color{Ink}Cigler's Hankel polynomials}
\fancyhead[R]{\small Research draft}
\fancyfoot[C]{\thepage}
\setlength{\parskip}{4pt}
\setlength{\parindent}{14pt}
\setlength{\emergencystretch}{2em}
\allowdisplaybreaks[2]
\setlist{nosep}

\title{\color{Ink}\LARGE A parity factorization and sharp stabilization\\[0.25em]
for Cigler's Hankel polynomials\\[0.5em]
\large A proposed proof of Conjecture 16 in arXiv:2111.14492v3}
\author{Research draft prepared for Vladimir Reshetnikov}
\date{20 September 2026}

\begin{document}
\maketitle
\thispagestyle{empty}
\begin{abstract}
We study the shifted Hankel determinants of Cigler's polynomial extension
$c_n(t)$ of the middle binomial coefficients. An even--odd decomposition
factors every normalized determinant into two principal minors of powers of
simple tridiagonal matrices. This gives a proof of polynomiality, strict
coefficient positivity, reciprocity, and the exact degree. An elementary
alternant evaluation of the two factors proves the stable-coefficient
assertions of Conjecture 16. More precisely, if
$a=\lfloor k/2\rfloor$, $b=\lfloor(k-1)/2\rfloor$, and
$\Ht_k(n;t)=(-1)^{k\binom n2}t^{-\binom n2}
\det(c_{k+i+j}(t))_{0\le i,j<n}$, then for $k\ge2$,
\[
\Ht_k(n;t)=\frac{1}{(1-t)^a(1-t^2)^{ab}}
-\binom{n+a}{a-1}t^{n+1}+O(t^{n+2}).
\]
Thus the conjectured stable range is optimal, and its first defect has a
closed formula. We also obtain fixed-size determinant evaluations, an
explicit nonminimal recurrence, and convergence on compact subsets of the
open unit disk. Exact integer and polynomial verification programs accompany
the article. A stronger rectangular-Schur identity suggested by computation
is separated from the proved results and is not used in any proof.
\end{abstract}

\medskip
\noindent\fcolorbox{Ink}{blue!3}{\parbox{0.94\linewidth}{\small
\textbf{Status and scope.} This is an AI-assisted research draft with a full
mathematical argument, not an independently refereed or proof-assistant-checked
result. The target is Conjecture 16 in Section 6 of Cigler's paper, not
Conjectures 13--15 in Section 5, which are already represented in the supplied
manifest. A targeted literature search did not locate a later solution of
this specific conjecture; that is not a guarantee of priority.}}

\clearpage
\tableofcontents
\clearpage

\section{The problem and the main result}
\label{sec:main}

\subsection{Why this problem is not a duplicate}
The supplied manifest \cite{manifest} catalogues a report entitled
\emph{Product formulas for ballot-polynomial Hankel determinants}, concerning
Conjectures 13--15 in Section 5 of Cigler's paper \cite{Cigler}.
Those conjectures concern the moment family denoted there by $a_n(t)$.
Here the family is $c_n(t)$ from equation (6), and the target is
\emph{Conjecture 16 in Section 6}, on pages 21--22 of version 3.
The relationship is methodological---polynomial moment sequences and their
Hankel determinants---but the input moments and the numbered problem differ.
We do not assume that any proof catalogued in the manifest is correct.

The source continues to present this statement as a conjecture in the
version located during the search. The arXiv record identifies version 3
as dated 30 December 2021 \cite{Cigler}. The bounded status audit in
Section~\ref{sec:audit} explains what was and was not established about
subsequent literature.

\subsection{Definitions and conventions}
A binomial coefficient $\binom rj$ is zero unless $0\le j\le r$.
Set $c_0(t)=1$, and for $r\ge1$ define
\begin{equation}
 c_r(t)=\sum_{j=1}^{r}
 \binom{\lfloor(r-1)/2\rfloor}{\lfloor(j-1)/2\rfloor}
 \binom{\lfloor r/2\rfloor}{\lfloor j/2\rfloor}t^{j-1}.
 \label{eq:source-moments}
\end{equation}
This is Cigler's equation (6), with its zero $j=0$ term omitted so that no
negative power of $t$ is displayed. For example,
\begin{align*}
 c_0&=1,&c_1&=1,&c_2&=1+t,\\
 c_3&=1+t+t^2,&
 c_4&=1+2t+2t^2+t^3,&
 c_5&=1+2t+4t^2+2t^3+t^4.
\end{align*}
For $k,n\ge0$, write
\begin{equation}
 D_k(n;t)=\det(c_{k+i+j}(t))_{0\le i,j<n},\qquad D_k(0;t)=1.
 \label{eq:D}
\end{equation}
For $k\ge1$ initially define, in $\Q(t)$,
\begin{equation}
 \Ht_k(n;t)=(-1)^{k\binom n2}\frac{D_k(n;t)}{t^{\binom n2}}.
 \label{eq:normalization}
\end{equation}
Polynomiality, including a well-defined value at $t=0$, will be proved.
It is important not to substitute $t=0$ into the quotient before this step.

Let
\begin{equation}
 a=\left\lfloor\frac k2\right\rfloor,\qquad
 b=\left\lfloor\frac{k-1}{2}\right\rfloor,\qquad
 \F_k(t)=\frac{1}{(1-t)^a(1-t^2)^{ab}}.
 \label{eq:F}
\end{equation}
An expression $O(t^r)$ in a formal-series identity means an element of
$t^r\Q[[t]]$; it has no analytic meaning unless one is explicitly stated.

\begin{theorem}[Conjecture 16, with a sharp first defect]
\label{thm:main}
For every $k\ge1$ and $n\ge0$, $\Ht_k(n;t)$ is a polynomial in $\Z[t]$ of
exact degree $(k-1)n$. Its constant and leading coefficients are both $1$,
every coefficient between them is strictly positive, and
\begin{equation}
 \Ht_k(n;t)=t^{(k-1)n}\Ht_k(n;t^{-1}).
 \label{eq:pal}
\end{equation}
For $k=1$, $\Ht_1(n;t)=1$. For $k\ge2$,
\begin{equation}
 \Ht_k(n;t)=\F_k(t)
 -\binom{n+a}{a-1}t^{n+1}+O(t^{n+2}).
 \label{eq:sharp}
\end{equation}
Consequently,
\begin{equation}
 [t^j]\Ht_k(n;t)=[t^j]\F_k(t)\quad(0\le j\le n),
 \label{eq:stable}
\end{equation}
and $n+1$ is the first degree of disagreement for every $k\ge2$.
\end{theorem}

For $k=2a$, equation~\eqref{eq:stable} has denominator
$(1-t)^a(1-t^2)^{a^2-a}$; for $k=2a+1$ it has denominator
$(1-t)^a(1-t^2)^{a^2}$. These are precisely the two stable-coefficient
assertions in the source. The first-defect formula \eqref{eq:sharp} is a
refinement, not part of the original statement.

\subsection{Proof architecture}
The proof has three independent ingredients. First, the relation
$c_{2r+2}=(1+t)c_{2r+1}$ turns a Hankel matrix into a two-by-two block matrix
whose determinant factors by row operations. Second, the resulting factors
are principal minors of powers of coefficientwise totally nonnegative
tridiagonal matrices. This establishes positivity and degree without any
alternant or Schur identity. Finally, small alternant determinants give the
first omitted terms of the two factors and hence the exact stable range.
All determinant identities needed below are proved in the article.

\section{Two auxiliary moment sequences}
\label{sec:moments}

Put $B=1+t^2$ and $s=1+t$, so that $s^2=B+2t$.
For a Laurent polynomial in $z$, $\CT_z$ denotes its constant coefficient.
Define
\begin{equation}
 Q(z)=B+t(z+z^{-1})=(1+tz)(1+t/z),
 \quad
 u_r=\CT_z(1+z)Q(z)^r,
 \quad
 v_r=\CT_z(1-z^2)Q(z)^r.
 \label{eq:uv-ct}
\end{equation}
Both $u_r$ and $v_r$ belong to $\Z[t]$.

\begin{lemma}[The pairing relation]
\label{lem:pair}
For every $r\ge0$,
\begin{equation}
 c_{2r+1}=u_r,\qquad c_{2r+2}=s u_r,\qquad
 s^2u_r-u_{r+1}=t v_r.
 \label{eq:pair}
\end{equation}
\end{lemma}
\begin{proof}
Expanding $Q(z)^r=(1+tz)^r(1+t/z)^r$ gives
\begin{equation}
 u_r=\sum_{j=0}^r\binom rj^2t^{2j}
 +\sum_{j=0}^{r-1}\binom rj\binom r{j+1}t^{2j+1}.
 \label{eq:u-binomial}
\end{equation}
Separating the even and odd indices of
\eqref{eq:source-moments} gives $c_{2r+1}=u_r$.
The coefficients of $t^{2j}$ and $t^{2j+1}$ in $(1+t)u_r$ are respectively
\[
 \binom rj\biggl(\binom rj+\binom r{j-1}\biggr),\qquad
 \binom rj\biggl(\binom rj+\binom r{j+1}\biggr).
\]
Pascal's identity identifies these with the coefficients of $c_{2r+2}$.
For the last assertion, observe that
\[
 (s^2-Q(z))(1+z)=t(1+z-z^2-z^{-1}).
\]
Since $Q(z)=Q(z^{-1})$, the constant terms of $zQ(z)^r$ and
$z^{-1}Q(z)^r$ agree. The unwanted two terms cancel, leaving $t v_r$.
\end{proof}

\subsection{Tridiagonal realizations}
Let $e_0,e_1,\ldots$ be the standard basis, let $S e_j=e_{j+1}$ be the
unilateral shift, and let $E_{00}=e_0e_0^{\mathsf T}$. For $\eps=0,1$ put
\begin{equation}
 M_\eps=B I+t(S+S^{\mathsf T})+\eps t E_{00}.
 \label{eq:M}
\end{equation}
Thus every off-diagonal entry adjacent to the diagonal is $t$, every
diagonal entry is $1+t^2$, and for $M_1$ the first diagonal entry has the
additional summand $t$. These are infinite banded matrices: every entry
of a fixed power and every minor used below is a finite calculation.

\begin{lemma}[Moment realization]
\label{lem:realization}
For $r\ge0$,
\begin{equation}
 u_r=(M_1^r)_{00},\qquad v_r=(M_0^r)_{00}.
 \label{eq:moment-M}
\end{equation}
\end{lemma}
\begin{proof}
For the first identity use Laurent polynomials
$f_h=z^h+z^{-h-1}$, $h\ge0$. Multiplication by $z+z^{-1}$ sends
$f_0$ to $f_0+f_1$ and sends $f_h$, $h>0$, to $f_{h-1}+f_{h+1}$.
Also $\CT f_h$ is $1$ for $h=0$ and $0$ otherwise. Multiplication by
$Q(z)$ in this basis therefore has matrix $M_1$, and
$\CT Q(z)^r f_0=(M_1^r)_{00}$. Replacing $z$ by $z^{-1}$ gives the
constant term defining $u_r$.

For the second identity use $g_h=z^h-z^{-h-2}$. Multiplication by
$z+z^{-1}$ sends $g_0$ to $g_1$ and, for $h>0$, sends $g_h$ to
$g_{h-1}+g_{h+1}$. Again $\CT g_h=\delta_{h0}$. This gives $M_0$ and
$\CT Q(z)^r(1-z^{-2})=v_r$.
\end{proof}

\section{Hankel determinants as principal minors}
\label{sec:gram}

For a matrix $A$, $A[m]$ means its principal submatrix on indices
$0,\ldots,m-1$, and $\det A[0]=1$. Define
\begin{equation}
 \U_a(m;t)=\det(M_1^a[m]),\qquad
 \V_a(m;t)=\det(M_0^a[m])\qquad(a,m\ge0).
 \label{eq:UV}
\end{equation}
Here $M_\eps^a[m]$ means $(M_\eps^a)[m]$, \emph{not}
$(M_\eps[m])^a$. Paths may leave the first $m$ indices and return before
the $a$ multiplications are complete.

\begin{lemma}[The auxiliary Gram identity]
\label{lem:gram}
For $a,m\ge0$,
\begin{align}
 \det(u_{a+i+j})_{0\le i,j<m}&=t^{m(m-1)}\U_a(m;t),\label{eq:gram-u}\\
 \det(v_{a+i+j})_{0\le i,j<m}&=t^{m(m-1)}\V_a(m;t).\label{eq:gram-v}
\end{align}
\end{lemma}
\begin{proof}
We prove both statements together. Write $M=M_\eps$ and define a moment
functional $L(X^r)=e_0^{\mathsf T}M^re_0$. Let
\[
 p_0=1,\quad p_1=X-B-\eps t,\quad
 p_{j+1}=(X-B)p_j-t^2p_{j-1}\quad(j\ge1).
\]
Induction gives $p_j(M)e_0=t^j e_j$. Consequently,
\begin{equation}
 L(p_i p_j)=t^{i+j}\delta_{ij},\qquad
 L(X^a p_i p_j)=t^{i+j}(M^a)_{ij}.
 \label{eq:orthogonality}
\end{equation}
The change from $1,X,\ldots,X^{m-1}$ to the monic polynomials
$p_0,\ldots,p_{m-1}$ is unit triangular. Changing both bases in the Hankel
Gram matrix therefore preserves its determinant. Taking the factors $t^i$
and $t^j$ out of its rows and columns gives
$t^{2\sum_{i=0}^{m-1}i}=t^{m(m-1)}$, as required.
This argument is an identity over $\Z[t]$ and does not require $t\ne0$.
\end{proof}

The normalized auxiliary quantities are thus polynomials before we use any
rational evaluation of them. In particular $\U_0(m)=\V_0(m)=1$ and
$\U_a(0)=\V_a(0)=1$.

\section{The parity factorization}
\label{sec:factor}

\begin{theorem}[Four parity cases]
\label{thm:factor}
For $a,m\ge0$,
\begin{align}
 \Ht_{2a+1}(2m;t)&=\U_a(m;t)\V_a(m;t),\label{eq:factor-oo}\\
 \Ht_{2a+1}(2m+1;t)&=\U_a(m+1;t)\V_a(m;t).\label{eq:factor-oe}
\end{align}
For $a\ge1$ and $m\ge0$,
\begin{align}
 \Ht_{2a}(2m;t)&=\U_a(m;t)\V_{a-1}(m;t),\label{eq:factor-eo}\\
 \Ht_{2a}(2m+1;t)&=(1+t)\U_{a-1}(m+1;t)\V_a(m;t).\label{eq:factor-ee}
\end{align}
In particular the normalization in \eqref{eq:normalization} is polynomial.
\end{theorem}
\begin{proof}
Reorder both rows and columns of the Hankel matrix by putting the even
indices first. Applying the same permutation on both sides contributes no
sign. Write $U_a(r)$ for the unnormalized matrix
$(u_{a+i+j})_{0\le i,j<r}$ and $V_a(r)$ for the analogous matrix of $v$'s.

For an odd shift $2a+1$, the four blocks, with the appropriate rectangular
sizes, are
\[
 \begin{pmatrix}
 (u_{a+i+j})& (s u_{a+i+j})\\
 (s u_{a+i+j})& (u_{a+i+j+1})
 \end{pmatrix}.
\]
Subtract $s$ times the corresponding even row from each odd row. The lower
left block becomes zero and the lower right block becomes
$(-t v_{a+i+j})$, by Lemma~\ref{lem:pair}. The upper left block has size
$m$ when $n=2m$ and size $m+1$ when $n=2m+1$. Therefore
\begin{align*}
 D_{2a+1}(2m)&=(-t)^m\det U_a(m)\det V_a(m),\\
 D_{2a+1}(2m+1)&=(-t)^m\det U_a(m+1)\det V_a(m).
\end{align*}
Using Lemma~\ref{lem:gram}, the total powers of $t$ are respectively
$2m^2-m=\binom{2m}{2}$ and $2m^2+m=\binom{2m+1}{2}$.
Also $(-1)^m=(-1)^{\binom n2}$ for both $n=2m$ and $n=2m+1$.
This proves the two odd-shift formulas.

For shift $2a$ and size $2m$, the reordered matrix is
\[
 \begin{pmatrix}sU_{a-1}(m)&U_a(m)\\ U_a(m)&sU_a(m)\end{pmatrix}.
\]
Work temporarily in $\Q(t)$ and subtract $s^{-1}$ times each bottom row
from the matching top row. The upper right block becomes zero, and the
upper left block is $(t/s)V_{a-1}(m)$. Hence
\[
 D_{2a}(2m)=t^m\det V_{a-1}(m)\det U_a(m).
\]
For size $2m+1$, the even block has $m+1$ rows. Subtract $s^{-1}$ times
even row $i+1$ from odd row $i$, for $0\le i<m$. The lower left block
vanishes; the lower right block becomes $(t/s)V_a(m)$, while the upper
left block is $sU_{a-1}(m+1)$. Thus
\[
 D_{2a}(2m+1)=s t^m\det U_{a-1}(m+1)\det V_a(m).
\]
The Gram identities give exactly the same two normalization exponents as
before. These are polynomial identities after denominators are cleared,
so the temporary use of $1/s$ causes no exception at $t=-1$.
\end{proof}

\section{Strict positivity, degree, and reciprocity}
\label{sec:positive}

For polynomials $P,Q$, write $P\coeffge Q$ if $P-Q$ has nonnegative
integer coefficients. A matrix is \emph{coefficientwise totally nonnegative}
if every finite minor, with row and column indices in increasing order,
is a polynomial with nonnegative coefficients.

\begin{lemma}[Total nonnegativity of the auxiliary matrices]
\label{lem:TP}
Both $M_0$ and $M_1$, and all of their nonnegative integer powers, are
coefficientwise totally nonnegative.
\end{lemma}
\begin{proof}
Since $S^{\mathsf T}S=I$ for the unilateral shift,
\begin{equation}
 M_0=(I+tS^{\mathsf T})(I+tS).
 \label{eq:bidiagonal}
\end{equation}
The two bidiagonal factors are totally nonnegative: a nonzero minor has
an order-preserving matching of its selected rows and columns and has a
nonnegative monomial weight. This can also be checked by expanding at the
first selected row; the bidiagonal zero pattern prevents a negatively
signed nonzero matching.

The Cauchy--Binet formula proves total nonnegativity of their product.
Although the matrices are infinite, each of these sums is finite because
the factors are banded. Next $M_1=M_0+tE_{00}$. A minor not selecting both
row $0$ and column $0$ is unchanged. A minor selecting both acquires
$t$ times the minor obtained by deleting its first row and first column.
The cofactor sign is positive because the selected indices are increasing.
Thus adding this corner entry preserves coefficientwise total
nonnegativity. Repeated Cauchy--Binet proves the assertion for powers.
\end{proof}

\begin{proposition}[Support and reciprocity of the factors]
\label{prop:factor-positive}
For $a,m\ge0$, $\U_a(m)$ and $\V_a(m)$ are reciprocal polynomials of
exact degree $2am$, with constant and leading coefficients $1$.
For $a,m\ge1$, every coefficient of $\U_a(m)$ is positive, whereas
$\V_a(m)$ has a positive coefficient at every even degree between $0$ and
$2am$ and no odd-degree terms. More quantitatively,
\begin{align}
 \U_a(m;t)&\coeffge(1+t+\cdots+t^{2m})^a,\label{eq:U-lower}\\
 \V_a(m;t)&\coeffge(1+t^2+\cdots+t^{2m})^a.\label{eq:V-lower}
\end{align}
\end{proposition}
\begin{proof}
The matrices satisfy $M_\eps(t^{-1})=t^{-2}M_\eps(t)$.
Taking the $a$th power and an $m$-row principal minor gives
\[
 \U_a(m;t)=t^{2am}\U_a(m;t^{-1}),\qquad
 \V_a(m;t)=t^{2am}\V_a(m;t^{-1}).
\]
At $t=0$, the matrices are $I$, so both determinants have constant term
$1$. Their exact degrees and leading coefficients follow.

The tridiagonal determinant recurrence, with $B=1+t^2$, gives
\[
 \det M_1[m]=1+t+\cdots+t^{2m},\qquad
 \det M_0[m]=1+t^2+\cdots+t^{2m}.
\]
For example the second sequence has initial values $1,B$ and recurrence
$P_m=BP_{m-1}-t^2P_{m-2}$; the first has initial values $1,B+t$
and the same recurrence. The displayed finite sums satisfy these recurrences.

Expand the principal minor of $M_\eps^a$ by repeated Cauchy--Binet. All
terms are coefficientwise nonnegative. Selecting the same index set
$\{0,\ldots,m-1\}$ at every intermediate stage supplies the term
$(\det M_\eps[m])^a$, proving the two lower bounds.
Finally $M_0(-t)=J M_0(t)J$, where $J=\diag(1,-1,1,-1,\ldots)$.
Its principal minors of powers are unchanged, so $\V_a(m)$ is even.
\end{proof}

\begin{proof}[Proof of the positivity and reciprocity part of Theorem~\ref{thm:main}]
Apply Theorem~\ref{thm:factor}. The degrees in its four lines are
\[
 4am,\quad 2a(2m+1),\quad 2m(2a-1),\quad (2a-1)(2m+1),
\]
which equal $(k-1)n$ in each case. All factors are reciprocal, including
$1+t$, and their constant terms are $1$.

For strict positivity, the full consecutive support of a nonconstant
$\U$ factor overlaps the translates arising from the even support of a
$\V$ factor. More explicitly, an interval $\{0,\ldots,L\}$ with $L\ge1$,
added to $\{0,2,\ldots,2R\}$, covers every integer from $0$ to $L+2R$.
The only exceptional-looking even-shift case is $a=1$, $n$ odd: there the
factor $1+t$ supplies precisely the needed interval of length one.
The cases $n=0$ and $k=1$ give the constant polynomial $1$ directly.
\end{proof}

\section{A small alternant evaluation}
\label{sec:alternants}

Positivity did not require symmetric-function identities. The stabilization
proof will also be elementary: the relevant formulas are quotients of
explicit alternating determinants. This section supplies their derivation,
including the classical moment determinant identity that is often used for
this purpose; compare the general framework in \cite{CK}.

\subsection{The Christoffel determinant identity}
\begin{lemma}[Moment modification]
\label{lem:christoffel}
Let $L$ be a moment functional over a characteristic-zero field, with monic
orthogonal polynomials $p_j$ of nonzero norms $h_j=L(p_j^2)$.
Put $H_m=\det(L(X^{i+j}))_{0\le i,j<m}$.
For independent parameters $\eta_1,\ldots,\eta_a$,
\begin{equation}
 \frac{\det\left(L\left(X^{i+j}\prod_{r=1}^a(X-\eta_r)\right)\right)_{0\le i,j<m}}{H_m}
 =(-1)^{am}
 \frac{\det(p_{m+j-1}(\eta_i))_{1\le i,j\le a}}
 {\Dlt(\eta)},
 \label{eq:christoffel}
\end{equation}
where $\Dlt(\eta)=\prod_{i<j}(\eta_j-\eta_i)$.
\end{lemma}
\begin{proof}
Form an $(m+a)$-square matrix whose first $m$ rows, indexed by $r$, are
$L(X^r X^j)$ for $0\le j<m+a$, and whose last $a$ rows are
$\eta_i^j$. Change the column polynomials from the monomials to
$p_0,\ldots,p_{m+a-1}$. The change is unit triangular. Orthogonality makes
the upper right block zero, and the determinant of the upper left block is
$\prod_{r=0}^{m-1}h_r=H_m$. Thus this determinant is
$H_m\det(p_{m+j-1}(\eta_i))$.

Alternatively change the column polynomials to
$1,X,\ldots,X^{a-1},P,XP,\ldots,X^{m-1}P$, where
$P=\prod_r(X-\eta_r)$. This is again a monic unit-triangular change.
Move the last $a$ rows to the top, contributing $(-1)^{am}$.
The top left block is a Vandermonde matrix and the top right block is zero.
The bottom right block is the modified moment matrix on the left of
\eqref{eq:christoffel}. Equating the two evaluations proves the identity.
\end{proof}

\subsection{The two explicit numerators}
For $a,m\ge0$ and $\delta\in\{0,1\}$, define
\begin{equation}
 N^{\delta}_{a,m}(t)=
 \det\left(
 \binom{j-1}{i-1}t^{j-i}
 -\binom{2m+2a-j+\delta}{i-1}
 t^{2m+2a-j+\delta-i+1}
 \right)_{1\le i,j\le a}.
 \label{eq:N}
\end{equation}
The first summand is interpreted as zero when $j<i$, without evaluating
its negative power. An empty determinant is $1$.

\begin{theorem}[Fixed-size evaluations]
\label{thm:alternant}
For all $a,m\ge0$,
\begin{align}
 \U_a(m;t)&=
 \frac{N^0_{a,m}(t)}{(1-t)^a(1-t^2)^{\binom a2}},\label{eq:Uclosed}\\
 \V_a(m;t)&=
 \frac{N^1_{a,m}(t)}{(1-t^2)^{\binom{a+1}{2}}}.
 \label{eq:Vclosed}
\end{align}
These identities hold in $\Q(t)$ and determine the polynomials already
constructed in \eqref{eq:UV}; apparent singularities at $t=1$ or $t=-1$
are removable. The determinant orders are $a$, independently of $m$.
\end{theorem}
\begin{proof}
Work first over $\Q(t)$. Let $C_0(y)=1$, $C_1(y)=y$, and
$C_{j+1}(y)=yC_j(y)-C_{j-1}(y)$; also put $C_{-1}=0$.
The orthogonal polynomials in Lemma~\ref{lem:gram} are
\[
 p_j^\eps(X)=t^j\left[C_j\left(\frac{X-B}{t}\right)
 -\eps C_{j-1}\left(\frac{X-B}{t}\right)\right].
\]
Define monic polynomials $A_j^\eps(y)=C_j(y)+\eps C_{j-1}(y)$ and put
\[
 \eta_i=B-t(x_i+x_i^{-1}).
\]
Then $p_j^\eps(\eta_i)=(-t)^jA_j^\eps(x_i+x_i^{-1})$.
The elementary recurrence gives
\begin{align}
 A_j^0(x+x^{-1})&=\frac{x^{j+1}-x^{-j-1}}{x-x^{-1}},\label{eq:A0}\\
 A_j^1(x+x^{-1})&=\frac{x^{j+1/2}-x^{-j-1/2}}
 {x^{1/2}-x^{-1/2}}=\sum_{r=-j}^{j}x^r.\label{eq:A1}
\end{align}
The half powers in the middle expression cancel; the expression on the
right is a Laurent polynomial.

Since the $A_j^\eps$ are monic in $y$, their first $a$ evaluation columns
have determinant $\Dlt(y)$, where $y_i=x_i+x_i^{-1}$.
Apply Lemma~\ref{lem:christoffel}, cancel the column powers of $-t$ and
the Vandermonde scale, and then specialize every $\eta_i$ to zero.
Choosing every $x_i=t$ gives
\begin{equation}
 \det(M_\eps^a[m])=
 t^{am}\left.
 \frac{\det(A_{m+j-1}^\eps(x_i+x_i^{-1}))}
 {\det(A_{j-1}^\eps(x_i+x_i^{-1}))}
 \right|_{x_1=\cdots=x_a=t}.
 \label{eq:character-step}
\end{equation}
The equal-parameter value is the limit of an alternating quotient, not a
quotient of two separately substituted zero determinants.

Reverse all columns in the numerator and denominator of this quotient and
use \eqref{eq:A0}--\eqref{eq:A1}. For $\delta=1-\eps$, multiplication
of row $i$ by the appropriate power of $x_i$ gives
\begin{equation}
 \left(\prod_{i=1}^a x_i\right)^m
 \frac{\det(A_{m+j-1}^\eps(x_i+x_i^{-1}))}
 {\det(A_{j-1}^\eps(x_i+x_i^{-1}))}
 =
 \frac{\det(x_i^{j-1}-x_i^{2m+2a-j+\delta})}
 {\det(x_i^{j-1}-x_i^{2a-j+\delta})}.
 \label{eq:alternant-ratio}
\end{equation}
At $x_i=t$, the prefactor on the left is exactly $t^{am}$.

The denominator on the right has the factorization
\begin{equation}
 \det(x_i^{j-1}-x_i^{2a-j+\delta})
 =\Dlt(x)\prod_{i=1}^a(1-x_i^{1+\delta})
 \prod_{i<j}(1-x_ix_j).
 \label{eq:den-factor}
\end{equation}
Here is a direct verification. Alternation supplies $\Dlt(x)$.
At $x_i=1$ a row vanishes, and for $\delta=1$ it also vanishes at $x_i=-1$.
At $x_j=x_i^{-1}$ the two rows are proportional, with proportionality
factor $-x_i^{-(2a-1+\delta)}$, so every $1-x_ix_j$ divides the determinant.
These factors have exactly its total degree. Its lowest homogeneous term
is $\Dlt(x)$, which fixes the multiplicative constant as $1$.

Finally, for polynomials $f_j$, Taylor expansion gives the confluent identity
\begin{equation}
 \left.\frac{\det(f_j(x_i))}{\Dlt(x)}\right|_{x_1=\cdots=x_a=t}
 =\det\left(\frac{f_j^{(i-1)}(t)}{(i-1)!}\right)_{1\le i,j\le a}.
 \label{eq:confluent}
\end{equation}
Indeed the lowest alternating term in the Taylor expansion around $t$ uses
the distinct exponents $0,1,\ldots,a-1$.
Apply \eqref{eq:confluent} to the numerator in
\eqref{eq:alternant-ratio}. It gives exactly \eqref{eq:N}.
The remaining factors in \eqref{eq:den-factor}, after equal specialization,
give the denominators in \eqref{eq:Uclosed} and \eqref{eq:Vclosed}.
\end{proof}

\begin{remark}
The quotients before equal specialization are familiar rectangular
orthogonal and symplectic character alternants. No character theory,
branching formula, or bounded Schur-sum identity is being imported: the
moment identity, the recurrence, the denominator factorization, and the
confluent limit above give a complete determinant proof.
\end{remark}

\section{Sharp stable coefficients}
\label{sec:stable}

\begin{lemma}[First nonconstant numerator term]
\label{lem:first-term}
For $a\ge1$, $m\ge0$, and $\delta\in\{0,1\}$, put
$r=2m+1+\delta$. Then
\begin{equation}
 N^\delta_{a,m}(t)=1-\binom{a+r-1}{a-1}t^r+O(t^{r+1}).
 \label{eq:N-first}
\end{equation}
\end{lemma}
\begin{proof}
Before equal specialization, expand
$\det(x_i^{j-1}-x_i^{2m+2a-j+\delta})$ by choosing either the low or the
high monomial in each column. If the high monomial is selected in column
$j$, the total degree increases, relative to $\Dlt(x)$, by
\[
 2m+2a-2j+1+\delta=r+2(a-j).
\]
The all-low choice is $\Dlt(x)$. The unique choice of smallest positive
additional degree is to select only the high monomial in column $a$;
its sign is negative and its additional degree is $r$.
Every other choice has strictly larger additional degree.

After division by $\Dlt(x)$ and equal specialization, the coefficient of
this first correction can be computed directly from
\eqref{eq:confluent}. The column exponents are
$0,1,\ldots,a-2,a-1+r$. The first $a-1$ columns are triangular with
unit diagonal, so the remaining determinant is
$\binom{a-1+r}{a-1}t^r$. This proves the assertion, including $a=1$.
\end{proof}

Define the formal series
\begin{equation}
 U_a^\infty(t)=\frac{1}{(1-t)^a(1-t^2)^{\binom a2}},\qquad
 V_a^\infty(t)=\frac{1}{(1-t^2)^{\binom{a+1}{2}}}.
 \label{eq:limits}
\end{equation}
Their constant terms are $1$. For $a=0$ both are exactly $1$.
Theorem~\ref{thm:alternant} and Lemma~\ref{lem:first-term} imply, for
$a\ge1$,
\begin{align}
 \U_a(m;t)&=U_a^\infty(t)
 -\binom{2m+a}{a-1}t^{2m+1}+O(t^{2m+2}),\label{eq:Udefect}\\
 \V_a(m;t)&=V_a^\infty(t)
 -\binom{2m+a+1}{a-1}t^{2m+2}+O(t^{2m+3}).\label{eq:Vdefect}
\end{align}
Multiplying by a series with constant term $1$ does not change the first
nonzero correction, which is why the same binomial coefficient appears in
these formulas.

\begin{proof}[Completion of Theorem~\ref{thm:main}]
For $n=2m$ the earliest correction in both
\eqref{eq:factor-oo} and \eqref{eq:factor-eo} comes from $\U_a(m)$,
at degree $2m+1=n+1$. Its magnitude is
$\binom{2m+a}{a-1}=\binom{n+a}{a-1}$.
The products of the corresponding limiting series are respectively
\[
 U_a^\infty V_a^\infty
 =\frac1{(1-t)^a(1-t^2)^{a^2}},\qquad
 U_a^\infty V_{a-1}^\infty
 =\frac1{(1-t)^a(1-t^2)^{a(a-1)}}.
\]
For $n=2m+1$ the earliest correction in both
\eqref{eq:factor-oe} and \eqref{eq:factor-ee} comes from $\V_a(m)$,
at degree $2m+2=n+1$, with magnitude
$\binom{2m+a+1}{a-1}=\binom{n+a}{a-1}$.
For the last parity case the limiting product simplifies as
\[
 (1+t)U_{a-1}^\infty V_a^\infty
 =\frac{1+t}{(1-t)^{a-1}(1-t^2)^{a^2-a+1}}
 =\frac1{(1-t)^a(1-t^2)^{a(a-1)}}.
\]
All multiplying factors have constant term $1$, so they do not alter the
first correction coefficient. When $a=1$, any factor with index $a-1=0$
is identically $1$, and the same argument applies. The positive binomial
coefficient proves optimality of the range. The case $k=1$ was already
obtained directly from the factorization.
\end{proof}

\begin{corollary}[Coefficient extraction without a Hankel determinant]
For $k\ge2$, $0\le j\le n$, and $a,b$ as in \eqref{eq:F},
\begin{equation}
 [t^j]\Ht_k(n;t)
 =\sum_{\ell=0}^{\lfloor j/2\rfloor}
 W(ab,\ell)W(a,j-2\ell),
 \label{eq:stable-sum}
\end{equation}
where $W(q,r)=\binom{q+r-1}{r}$ for $q\ge1$, and
$W(0,r)=1$ for $r=0$ and $0$ otherwise.
\end{corollary}
\begin{proof}
Expand $(1-t)^{-a}$ and $(1-t^2)^{-ab}$ and use
\eqref{eq:stable}. The convention for $W(0,r)$ handles $k=2$ without
an artificial negative upper binomial index.
\end{proof}

\section{Examples and checks on the statement}
\label{sec:examples}

\subsection{The first three shifts}
The factorization immediately gives
\begin{equation}
 \Ht_1(n;t)=1,\qquad
 \Ht_2(n;t)=1+t+\cdots+t^n.
 \label{eq:k12}
\end{equation}
For $k=3$, insert
$\U_1(m)=1+t+\cdots+t^{2m}$ and
$\V_1(m)=1+t^2+\cdots+t^{2m}$ in the two parity cases. Both simplify to
\begin{equation}
 \Ht_3(n;t)=\frac{(1-t^{n+1})(1-t^{n+2})}{(1-t)(1-t^2)}.
 \label{eq:k3}
\end{equation}
These small-shift formulas are already printed in Cigler's paper; they
serve as indexing and sign checks, not as new results.

\subsection{A first-defect example}
For $k=4$, the limiting series is
\[
 \F_4(t)=\frac1{(1-t)^2(1-t^2)^2}
 =1+2t+5t^2+8t^3+14t^4+\cdots.
\]
At $n=2$ the determinant gives
\[
 \Ht_4(2;t)=1+2t+5t^2+4t^3+5t^4+2t^5+t^6.
\]
The first three coefficients agree. At degree $3=n+1$ the deficit is
$8-4=4=\binom{2+2}{2-1}$, exactly as predicted.
At $n=3$,
\[
 \Ht_4(3;t)=1+2t+5t^2+8t^3+9t^4+9t^5+8t^6+5t^7+2t^8+t^9,
\]
and the first deficit is $14-9=5=\binom{3+2}{2-1}$.

These examples also delimit the result. The coefficients $5,4,5$ in
$\Ht_4(2;t)$ have a strict valley. Thus positivity and reciprocity do
\emph{not} imply unimodality here. No unimodality or log-concavity assertion
is included in Theorem~\ref{thm:main}.

\subsection{Signs and exceptional parameters}
For real $t>0$, strict coefficient positivity gives the exact sign
\begin{equation}
 \sgn D_k(n;t)=(-1)^{k\binom n2}.
 \label{eq:sign}
\end{equation}
At $t=0$ the normalized polynomial has value $1$, although the original
Hankel determinant vanishes when $n\ge2$. At $t=1$ the source moments
are the middle binomial coefficients. At $t=-1$ the factor $1+t$ in
\eqref{eq:factor-ee} shows that every positive even-shift, odd-size determinant
vanishes. These statements follow from polynomial identities; no argument
uses an undefined specialization of a rational formula.

\section{Algorithms and further consequences}
\label{sec:consequences}

\subsection{Exact computation at fixed shift}
Theorem~\ref{thm:alternant} evaluates the factors through determinants of
order at most $\lfloor k/2\rfloor$, even when $n$ is very large.
A practical algorithm is to determine the parity of $k,n$, compute the
one or two numerators in \eqref{eq:N}, divide by their proved exact
denominators, and multiply according to Theorem~\ref{thm:factor}.
For $t=\pm1$, use the principal-minor definition or take the polynomial
limit instead of evaluating a zero denominator.

The reduction is in determinant \emph{order}, not in bit complexity.
The entries contain powers whose exponents grow with $n$, and an expanded
answer has degree $(k-1)n$. Once entries are constructed, ordinary dense
elimination uses a cubic number of field operations in the small determinant
order; large integers and polynomial output still carry their usual costs.
For only the first $j\le n$ coefficients, \eqref{eq:stable-sum} avoids
both determinant constructions.

\subsection{A concrete, nonminimal recurrence}
\begin{proposition}[Rationality at each fixed shift]
\label{prop:recurrence}
For $k\ge2$, set
\[
 d_k=\begin{cases}(a-1)^2,&k=2a,\\ a(a-1),&k=2a+1.\end{cases}
\]
The series $\sum_{n\ge0}\Ht_k(n;t)x^n$ is rational over $\Q(t)$,
with a denominator dividing
\begin{equation}
 \prod_{q=0}^{k-1}(1-t^{2q}x^2)^{d_k+1}.
 \label{eq:coarse-den}
\end{equation}
Equivalently, with $E$ denoting forward shift in $n$,
\begin{equation}
 \prod_{q=0}^{k-1}(E^2-t^{2q})^{d_k+1}\Ht_k(n;t)=0.
 \label{eq:coarse-rec}
\end{equation}
The bound is not asserted to be minimal and is not a resolution of all the
more precise denominator and numerator assertions in Conjecture 18.
\end{proposition}
\begin{proof}
Expand \eqref{eq:N} column by column. A term selecting $q$ high columns
has the form $P(m;t)t^{2qm}$, where $P$ is a polynomial in $m$ with
coefficients in $\Q(t)$. Its degree in $m$ is at most
$\sum_{i=1}^a(i-1)=\binom a2$, because the binomial in row $i$ has degree
at most $i-1$. The denominators in Theorem~\ref{thm:alternant} do not
depend on $m$. Thus each auxiliary factor is a finite sum of
$m^j(t^{2q})^m$, with $0\le q\le a$ and $j\le\binom a2$.

The products in the two parity subsequences have $0\le q\le k-1$.
Their degree bounds are $2\binom a2=a(a-1)$ in the odd-shift case and
$\binom a2+\binom{a-1}{2}=(a-1)^2$ in the even-shift case.
Replacing $m$ by $m+1$ changes neither bound. A polynomial of degree at
most $d$ times $\lambda^m$ is annihilated by $(E_m-\lambda)^{d+1}$,
as follows by taking successive finite differences. Interleaving the
even and odd subsequences replaces $E_m$ by $E^2$, proving the claim.
\end{proof}

For the unadjusted normalized determinant
$d_k(n;t)=t^{-\binom n2}D_k(n;t)$, the even-$k$ recurrence is unchanged.
For odd $k$, the period-four sign replaces $1-t^{2q}x^2$ in the displayed
coarse denominator by $1+t^{2q}x^2$.

\subsection{Actual convergence, not just formal stabilization}
\begin{proposition}[Uniform convergence inside the unit disk]
\label{prop:analytic}
Fix $k\ge2$ and $0<r<1$. There is a constant $C_{k,r}$ such that, for all
$n\ge0$,
\begin{equation}
 \sup_{|t|\le r}|\Ht_k(n;t)-\F_k(t)|
 \le C_{k,r}(n+1)^{\binom a2}r^{n+1}.
 \label{eq:analytic}
\end{equation}
In particular the polynomials converge to $\F_k$ locally uniformly on
$|t|<1$.
\end{proposition}
\begin{proof}
In the expansion of $N^\delta_{a,m}-1$, every nonzero term has degree at
least $2m+1+\delta$ in $t$. There are finitely many determinant terms for
fixed $a$, and each coefficient is bounded by a constant times
$(m+1)^{\binom a2}$, by the row-degree estimate just used.
It follows, uniformly on $|t|\le r$, that
\[
 N^\delta_{a,m}(t)-1
 =O_{a,r}\bigl((m+1)^{\binom a2}r^{2m+1+\delta}\bigr).
\]
The two fixed denominators in Theorem~\ref{thm:alternant} are bounded
away from zero on this disk. Hence the factors converge uniformly to
\eqref{eq:limits}, with these bounds, and in particular remain uniformly
bounded as $m$ varies. Subtract the limiting products in the four parity
formulas. The earliest error has exponent $n+1$, exactly as in
Section~\ref{sec:stable}; the largest polynomial exponent is $\binom a2$.
The finitely many small cases can be absorbed in $C_{k,r}$.
\end{proof}

\section{A stronger identity suggested by the computations}
\label{sec:schur}

The discovery phase suggested a compact formula not needed for any theorem
above. We record it as a separate conjecture so that a promising observation
is not silently promoted to a proof.

For an alphabet $X=(x_1,\ldots,x_k)$, define $h_j(X)$ by
$\sum_{j\ge0}h_j(X)z^j=\prod_{i=1}^k(1-x_i z)^{-1}$, with $h_j=0$ for
$j<0$. The rectangular Schur polynomial is
\[
 s_{(n^a)}(X)=\det(h_{n-i+j}(X))_{1\le i,j\le a}.
\]

\begin{conjecture}[Experimental rectangular-Schur formula]
\label{conj:schur}
With $a=\lfloor k/2\rfloor$ and $b=\lfloor(k-1)/2\rfloor$,
\begin{equation}
 \Ht_k(n;t)=s_{(n^a)}
 (\underbrace{1,\ldots,1}_{a},\ t,\
 \underbrace{t^2,\ldots,t^2}_{b}).
 \label{eq:schur-candidate}
\end{equation}
\end{conjecture}

The conjecture agrees with all 1,248 exact evaluations in the range recorded
in Section~\ref{sec:verify}. For instance, at $k=4,n=2$, its right-hand
side is $h_2^2-h_1h_3$ in the alphabet $(1,1,t,t^2)$, giving the
polynomial displayed in Section~\ref{sec:examples}.

A proof of \eqref{eq:schur-candidate} would give an attractive unified
symmetric-function interpretation. It may be approachable through
factorizations of rectangular Schur polynomials into classical-character
alternants. This article does not assert such a factorization without
proof. The proof of Conjecture 16 instead uses the parity factors,
coefficientwise total nonnegativity, and the fully derived alternant
formulas. The larger conjecture can be deleted from this article without
changing any proved statement.

\section{Reproducibility and proof audit}
\label{sec:audit}

\subsection{What was checked computationally}
\label{sec:verify}
The accompanying program \texttt{code/verify.py} uses exact integers and
exact polynomials in $\Z[t]$. It does not infer symbolic identities from
floating-point samples or from interpolation. The direct moment routine
implements the source definition \eqref{eq:source-moments}; the principal
minor routine independently multiplies the banded matrices; and the small
determinant routine implements \eqref{eq:N}.

\begin{center}
\small
\begin{tabular}{@{}p{0.35\linewidth}p{0.48\linewidth}r@{}}
\toprule
Check & Range & Cases\\
\midrule
Original Hankel versus parity formula &
$1\le k\le12$, $0\le n\le12$, eight integer parameters & 1,248\\[3pt]
Auxiliary corner versus alternant &
$0\le a\le7$, $0\le m\le8$, both factors, six parameters & 864\\[3pt]
Three-way polynomial equality; degree, positivity, reciprocity,
stability, first defect & $1\le k\le10$, $0\le n\le8$, in $\Z[t]$ & 90\\[3pt]
Auxiliary first numerator term &
$1\le a\le6$, $0\le m\le6$, both numerators, in $\Z[t]$ & 84\\[3pt]
Experimental Schur formula only &
$1\le k\le12$, $0\le n\le12$, eight integer parameters & 1,248\\
\bottomrule
\end{tabular}
\end{center}

The eight parameters are $-3,-2,-1,0,1,2,3,5$. The six-parameter
alternant test omits $\pm1$, where its displayed denominator vanishes;
these parameters are nevertheless covered by the original-versus-corner
checks. All checks passed with Python 3.13.5 and SymPy 1.14.0.
The execution record is in \texttt{data/verification.json}, with the
console transcript in \texttt{data/verification.log} and coefficient
arrays in \texttt{data/sample\_polynomials.json}. A supplementary script
checks the specialization products in Appendix~\ref{app:one} in 234 cases
and the recurrence in Proposition~\ref{prop:recurrence} in 126 cases; its
exact ranges and results are recorded in
\texttt{data/additional\_verification.json}.
These finite checks supplement the all-index proofs; they are not a
substitute for them.

\subsection{Logical dependencies and possible failure points}
The main theorem depends only on the binomial definition, the Laurent
basis calculation, elementary block row operations, Cauchy--Binet,
monic orthogonal-polynomial recurrences, and alternating determinants.
The proof of the moment-modification identity is included rather than
used as an unexplained import. The exact role of each delicate point is
as follows.

The odd/even row reorder is applied to both rows and columns, so there is
no permutation sign. The only sign remaining is $(-1)^{\lfloor n/2\rfloor}$
for odd shifts, which equals the prescribed $(-1)^{\binom n2}$.
The temporary division by $1+t$ in the even-shift factorization is removed
before specialization. The auxiliary normalization is justified by an
identity in $\Z[t]$, not by dividing a determinant whose vanishing has
not been established. Equal-variable alternants are evaluated after
cancelling the Vandermonde, not by a $0/0$ substitution. Finally,
strict positivity is obtained from explicit Cauchy--Binet summands;
nonnegativity alone would not have been sufficient.

\subsection{Literature status and exclusions}
The source pages and their mathematical notation were inspected directly.
Searches by the paper identifier, its title, and the exact numbered
conjecture did not locate a later proof of this specific statement.
A search hit for Manuel Kauers's \emph{A Proof of Conjecture 16} \cite{KauersIndex} was
investigated separately: indexed material identifies its target as the
sequence A339987, not Cigler's $c_n(t)$ conjecture. It was therefore not
counted as a solution of the present problem. The full talk PDF was not
successfully fetched, so this disambiguation rests on the indexed author
material rather than a complete reading of the talk.

The article makes no absolute priority claim. It also does not claim to
settle all parts of Cigler's Conjectures 17 and 18, or the separate Schur
formula \eqref{eq:schur-candidate}. What is claimed mathematically is the
full argument for Conjecture 16 and the additional results explicitly
proved here. Independent expert checking remains appropriate before
publication or addition to a collection of established theorems.

\appendix
\section{The removable specialization at \texorpdfstring{$t=1$}{t=1}}
\label{app:one}

The factors have convenient product evaluations at $1$:
\begin{align}
 \U_a(m;1)&=
 \prod_{r=1}^a\frac{2m+2r-1}{2r-1}
 \prod_{1\le r<s\le a}\frac{2m+r+s-1}{r+s-1},\label{eq:Uone}\\
 \V_a(m;1)&=
 \prod_{r=1}^a\frac{m+r}{r}
 \prod_{1\le r<s\le a}\frac{2m+r+s}{r+s}.
 \label{eq:Vone}
\end{align}
Empty products have value $1$.
To see these formulas directly from the proved alternants, set
$x_i=e^{z_i}$ before taking the equal-parameter limit in
\eqref{eq:character-step}. For $\V$, the numerator contains
$\sinh((m+r)z_i)$ and the denominator $\sinh(rz_i)$, with $1\le r\le a$.
For $\U$, replace $r$ by $r-\tfrac12$.
The lowest nonzero alternating term of
$\det(\sinh(\rho_j z_i))$ is a common nonzero constant times
\[
 \left(\prod_i z_i\right)\Dlt(z_1^2,\ldots,z_a^2)
 \left(\prod_j\rho_j\right)\Dlt(\rho_1^2,\ldots,\rho_a^2).
\]
This follows by selecting the powers $1,3,\ldots,2a-1$ in the Taylor
series. In the ratio, the differences between the $\rho_j$ cancel,
whereas the sums give the second products in
\eqref{eq:Uone}--\eqref{eq:Vone}. These specialize the new parity formula
to the classical middle-binomial Hankel setting already treated in the
source; the unweighted setting itself is not claimed as new.

\clearpage
\section{How to run the package}

The archive contains the article source and PDF, three Python files, exact
verification output, example coefficient arrays, and a status/source audit.
The commands, from the archive's root directory, are
\begin{verbatim}
python -m pip install -r requirements.txt
python code/verify.py
python code/verify_additional.py
pdflatex -interaction=nonstopmode -halt-on-error article.tex
pdflatex -interaction=nonstopmode -halt-on-error article.tex
\end{verbatim}
The \texttt{Makefile} provides \texttt{make verify} and \texttt{make pdf}.
The verification program stops at the first failed assertion or inexact
division. It records the mathematical test ranges and software versions,
not just a generic success message. Its stronger Schur checks are labelled
\texttt{EXPERIMENTAL} in both the code and the output.

\begin{thebibliography}{9}
\bibitem{Cigler}
Johann Cigler,
\emph{Hankel determinants of middle binomial coefficients and conjectures
for some polynomial extensions and modifications},
arXiv:2111.14492v3, 30 December 2021.
The present target is Conjecture 16, pp.~21--22; the input moments are
equation (6), p.~3.
\url{https://arxiv.org/abs/2111.14492v3}.

\bibitem{CK}
Johann Cigler and Christian Krattenthaler,
\emph{Hankel determinants of linear combinations of moments of orthogonal
polynomials}, arXiv:2003.01676, 2020.
\url{https://arxiv.org/abs/2003.01676}.

\bibitem{manifest}
Vladimir Reshetnikov,
\emph{A manifest of docs/reports: Seventy-one independent mathematical
research packages, classified by subject}, supplied file
\texttt{manifest.tex}, dated 20 September 2026.
See the entry \emph{Product formulas for ballot-polynomial Hankel
determinants} under \emph{Catalan and ballot moments}.
The manifest records report claims, not independent proof verification.
\bibitem{KauersIndex}
Manuel Kauers,
\emph{A Proof of Conjecture 16}, author-hosted lecture material,
\url{https://www.algebra.uni-linz.ac.at/people/mkauers/publications/kauers26a.pdf}.
Search-index excerpts consulted 20 September 2026 identify the target as
A339987. The complete PDF was not successfully retrieved; this reference
is used only to disambiguate the numbered problem, not as a mathematical
proof source.
\end{thebibliography}
\end{document}
